\section{대수학 문제 (학사 수준, 신기하고 재밌는 것)}

2026-10-10. 커뮤니티 배포용. 난이도: 1 \textless{} 2 \textless{} 3 ≈ 4
\textless{} 5. 5번은 문제만 먼저 던지고 답을 접어두면 반응이 좋다.

\subsection{1. 불 환은 가환}

환 \(R\)의 모든 원소가 \(x^2=x\)를 만족하면 \(R\)은 가환환이다.

\textbf{풀이.} \((x+x)^2=x+x\)에서 \(4x=2x\), 즉 \(2x=0\)이고 \(x=-x\).
\((x+y)^2=x+y\)를 전개하면 \(xy+yx=0\), 따라서 \(xy=-yx=yx\).
\(\blacksquare\)

\subsection{2. 유한체의 모든 원소는 두 제곱수의 합}

\(q\)가 홀수인 소수 거듭제곱이면 \(\mathbb F_q\)의 모든 원소는 두
제곱수의 합이다. 이로부터 모든 소수 \(p\)에 대해
\(x^2+y^2+1\equiv0\pmod p\)가 해를 가짐을 보여라.

\textbf{풀이.} 제곱수 집합 \(S=\{x^2\}\)는 크기 \(\tfrac{q+1}{2}\)이고,
\(T=\{c-y^2\}\)도 크기 \(\tfrac{q+1}{2}\)이다. \(|S|+|T|=q+1>q\)이므로
비둘기집 원리로 \(S\cap T\ne\varnothing\), 즉 \(c=x^2+y^2\). \(c=-1\)을
넣으면 따름정리. \(p=2\)는 \((1,0)\). \(\blacksquare\)

\subsection{3. 정육면체 칠하기 (번사이드)}

여섯 면을 \(n\)색으로 칠할 때 회전으로 겹치는 것을 같게 보면 몇
가지인가?

\begin{longtable}[]{@{}llll@{}}
\toprule\noalign{}
회전 & 개수 & 면의 순환 구조 & 고정되는 칠하기 \\
\midrule\noalign{}
\endhead
\bottomrule\noalign{}
\endlastfoot
항등원 & 1 & \(1^6\) & \(n^6\) \\
면 축 \(\pm90^\circ\) & 6 & \(1^24^1\) & \(n^3\) \\
면 축 \(180^\circ\) & 3 & \(1^22^2\) & \(n^4\) \\
꼭짓점 축 \(\pm120^\circ\) & 8 & \(3^2\) & \(n^2\) \\
모서리 축 \(180^\circ\) & 6 & \(2^3\) & \(n^3\) \\
\end{longtable}

\[
\frac{n^6+3n^4+12n^3+8n^2}{24}
\]

\(n=2\)이면 10, \(n=3\)이면 57.

\subsection{\texorpdfstring{4. \(60^\circ\)는 삼등분할 수 없다}{4. 60\^{}\textbackslash circ는 삼등분할 수 없다}}

\(\cos3\theta=4\cos^3\theta-3\cos\theta\)에 \(\theta=20^\circ\),
\(y=2\cos20^\circ\)를 넣으면

\[
y^3-3y-1=0
\]

유리근 후보 \(\pm1\)이 근이 아니므로 기약이고
\([\mathbb Q(\cos20^\circ):\mathbb Q]=3\). 작도 가능한 수의 차수는 2의
거듭제곱이므로 불가능. \(\blacksquare\)

\subsection{5. 기약인데 모든 소수에서 가약}

\(x^4+1\)은 \(\mathbb Z[x]\)에서 기약이지만 모든 소수 \(p\)에 대해
\(\mathbb F_p[x]\)에서 가약이다.

\textbf{풀이.} \(x\to x+1\)을 넣은 \(x^4+4x^3+6x^2+4x+2\)에
아이젠슈타인(\(p=2\)). \(p=2\)이면 \(x^4+1=(x+1)^4\). 홀수 \(p\)에서는
비제곱수 둘의 곱이 제곱수이므로 \(-1,\,2,\,-2\) 중 하나는 제곱수다.

\begin{itemize}
\tightlist
\item
  \(-1=a^2\): \((x^2-a)(x^2+a)\)
\item
  \(2=a^2\): \((x^2+ax+1)(x^2-ax+1)\)
\item
  \(-2=a^2\): \((x^2-ax-1)(x^2+ax-1)\) \(\blacksquare\)
\end{itemize}

\textbf{배경.} 분해체 \(\mathbb Q(\zeta_8)\)의 갈루아 군은
\(\mathbb Z/2\times\mathbb Z/2\)라서 4개의 근을 한 번에 도는 4-순환
원소가 없다. 체보타레프 밀도 정리에 의해, \(n\)차 기약다항식이 모든
소수에서 가약일 필요충분조건은 갈루아 군이 근 위에서 \(n\)-순환을 갖지
않는 것이다.
